

\documentclass[12pt,a4paper]{article}
%数学宏包
\usepackage{amsmath}
\usepackage{amsbsy}
\usepackage{mathtools}
\usepackage{amssymb}
\usepackage{amsthm}
\usepackage{thmtools}
%图片
\usepackage{ graphicx}
%中文
\usepackage{ctex}
%页面设置
\usepackage{geometry}
\geometry{left=2cm,right=2cm,top=2cm,bottom=2cm}
% 标题格式设置
\usepackage{titlesec}
\usepackage{titletoc}

    
\renewcommand{\(}{\left(}
\renewcommand{\)}{\right)}
\newcommand{\D}{\Delta}

\renewcommand{\l}{\left| }
\renewcommand{\r}{\right|}
  \newcommand{\bvec}[1]{\overrightarrow{#1}}
\title{有关距离函数的思考}
    
\begin{document}
\maketitle

\section{从椭圆引入}

对椭圆C:
$$ \frac{x^2}{a^2}+\frac{y^2}{b^2}=1$$
我们可以通过许多普通方式求出C在 $\(x_0,y_0\)$的法向量
$$ \vec{n}=\( \frac{x_0}{a^2},\frac{y_0}{b^2} \)$$
下面我们从另一个角度来看：
    
我们将$C: \frac{x^2}{a^2}+\frac{y^2}{b^2}=1 $看作函数$z=f\(x,y\)=\frac{x^2}{a^2}+\frac{y^2}{b^2}$的等值线
    
由全微分方程$$dz=f'_x \Delta x+f'_y \Delta y$$
当$\sqrt{{\D x}^2+{\D y}^2} \to 0$时,$\D z=dz$.{ }
当$\( x,y \)$ 于$C$上运动时，$z=1$,
$\D z=0$
    
即$$\D z=f'_x \Delta x+f'_y \Delta y=0$$
或写成在$ \(x_0,y_0 \)$的形式
\begin{align*}
\D z  =&   f \(x_0+\D x,y_0+\D y \) -f \(x_0,y_0 \)    \\
        =&  f'_x \(x_0,y_0 \)\Delta x+f'_y\(x_0,y_0 \) \Delta y\\
        =&0
\end{align*}

设
\begin{align*}
\begin{cases}
\vec{v}     =   \(\D x,\D y\) \\
\vec{n}     =   \( f'_x \(x_0,y_0 \),f'_y\(x_0,y_0 \) \)
\end{cases}
\end{align*}

则$$\vec{v}  \cdot  \vec{n} =0$$

则$C$在$  \( x_0 , y_0 \)  $处的法向量为
$$   \vec{n}     =   \( f'_x \(x_0,y_0 \),f'_y\(x_0,y_0 \) \)$$

对于函数$z=f\(x,y\)=\frac{x^2}{a^2}+\frac{y^2}{b^2}$
\begin{align*}
\begin{cases}
 f'_x \(x_0,y_0 \)=&   2 \frac{x_0}{a^2}  \\
f'_y\(x_0,y_0 \)=&    2\frac{y_0}{b^2} 
\end{cases}
\end{align*}
得$$ \vec{n}=\(2  \frac{x_0}{a^2},2  \frac{y_0}{b^2} \)$$

\section{如果从源头出发}

从椭圆的定义出发，$PF_1+PF_2=2a$，椭圆方程可写成
$$ C:  \sqrt{\( x+c \)^2 +y^2}+ \sqrt{ \( x-c \)^2 +y^2 }=2a  $$
该方程显然与$C$一点儿区别没有.设函数
$$ z= g \( x,y \) =\sqrt{\( x+c \)^2 +y^2}+ \sqrt{ \( x-c \)^2 +y^2 } $$
则$C$是一条等值线 $z=2a$.
再利用上述方法，

\begin{align*}
\begin{cases}
g'_x \( x_0,y_0 \)=\frac{   x_0+c}{  \sqrt{\( x_0+c \)^2 +y_0^2}} + \frac{   x_0-c}{  \sqrt{\( x_0-c \)^2 +y_0^2}}            \\
g'_y  \( x_0,y_0 \)=\frac{   y_0}{  \sqrt{\( x_0+c \)^2 +y_0^2}} + \frac{   y_0}{  \sqrt{\( x_0-c \)^2 +y_0^2}}    
\end{cases}
\end{align*}

则$C$在$  \( x_0 , y_0 \)  $处的法向量为
\begin{align*}
  \vec{n}    & =   \( g'_x \(x_0,y_0 \),g'_y\(x_0,y_0 \) \)     \\
 \vec{n}     &=   \( \frac{   x_0+c}{  \sqrt{\( x_0+c \)^2 +y_0^2}} + \frac{   x_0-c}{  \sqrt{\( x_0-c \)^2 +y_0^2}},     \frac{   y_0}{  \sqrt{\( x_0+c \)^2 +y_0^2}} + \frac{   y_0}{  \sqrt{\( x_0-c \)^2 +y_0^2}}\)    
\end{align*}
得
$$     \vec{n}    =    \frac{  \bvec{F_1P}  }{   \l   F_1P  \r  }   +\frac{  \bvec{  F_2P  }  }{   \l F_2P  \r  }  $$

\section{   一点启发 }
由上一部分得：
设$P{ }： \( x,y \) $,与$P$有关的函数$ z= g \( x,y \)= F_1P +F_2P $






\end{document}
